\section{Boundaries}

Digit and Benford-style tests are weak forensic screens. They are easy to run,
easy to rank, and easy to misread. A high score can justify asking why a count
table has a particular digit pattern; it cannot answer that question by itself.
It is not proof of fraud, not a certification failure, and not a substitute for
risk-limiting audits or ordinary canvass reconciliation.

The main technical boundary is the Benford assumption. Many election-count
tables do not span the orders of magnitude needed for leading-digit behavior.
Precincts are often intentionally similar in size; contests impose upper bounds;
turnout, registration, and reporting-unit rules shape the available numbers.
Second-digit tests add another layer of fragility and should be treated with
particular caution. Last-digit tests avoid the Benford premise but remain
sensitive to aggregation, rounding, transcription, and reporting conventions.

Accordingly, W.04 reports should never be used alone. They should be a
low-weight signal in a broader investigation queue, read alongside source
hashes, lineage, audit records, batch information, geography, and human review.
The lead boundary from W.01 remains controlling: W-series outputs prioritize
investigation; they do not prove fraud, certify correctness, or replace audits.
